\chapter*{CONCLUSIONES PRELIMINARES}
\addcontentsline{toc}{chapter}{CONCLUSIONES PRELIMINARES}
\chaptermark{CONCLUSIONES PRELIMINARES}

\begin{enumerate}
  \item Se identificó y categorizó un conjunto de 10 criterios críticos para
        la evaluación de tecnologías de interfaz de usuario en el sector
        bancario, organizados en dos dimensiones: técnica (60\% del peso
        total: seguridad, rendimiento, mantenibilidad, calidad del código
        generado por IA, accesibilidad y compatibilidad) y regulatoria
        (40\%: cumplimiento normativo, integración con sistemas legacy,
        alta disponibilidad y gobernanza del uso de IA generativa). A
        diferencia de una primera versión de este modelo, los factores
        organizacionales (ecosistema, talento, alineación estratégica) se
        excluyeron deliberadamente de la ponderación, por las razones
        expuestas en la conclusión 6. Esta reorganización confirma que las
        metodologías genéricas de evaluación tecnológica son
        estructuralmente insuficientes para el sector financiero al no
        contemplar con la especificidad requerida ni las exigencias
        regulatorias ni los criterios emergentes asociados a la IA
        generativa.

  \item Los instrumentos metodológicos diseñados --matriz de ponderación AHP
        con validación mediante Técnica Delphi, scorecard técnico comparativo
        y rúbricas de evaluación cuantitativa-- demostraron consistencia
        inter-evaluador con razón de consistencia $CR = 0.06$. La
        objetividad del proceso se incrementó significativamente respecto a
        evaluaciones previas ad hoc de la institución estudiada, eliminando
        la variabilidad subjetiva entre departamentos documentada en el
        diagnóstico inicial y garantizando la trazabilidad de cada decisión
        de ponderación.

  \item El proceso estructurado de toma de decisiones de cuatro fases
        --rediseñado para centrarse exclusivamente en las dimensiones
        técnica y regulatoria, con la Fase 3 redefinida como
        ``Evaluación Regulatoria y de Gobernanza de IA''-- integró
        efectivamente ambas dimensiones en un flujo decisional coherente,
        sin requerir un rol dedicado a la validación estratégica de
        negocio. La participación coordinada de arquitectos empresariales,
        equipos de seguridad y oficiales de cumplimiento redujo los
        patrones de fragmentación decisional identificados en la revisión
        de literatura.

  \item La validación empírica en el caso de estudio --una institución
        bancaria mediana peruana con 20 profesionales de desarrollo
        organizados bajo metodología Scrum, con adopción incipiente de
        asistentes de código IA en el 60\% del equipo, y cerca de medio
        millón de cuentas activas-- confirmó la aplicabilidad práctica de la
        metodología propuesta. Angular obtuvo la puntuación ponderada más
        alta ($S = 8.84$) sobre React ($S = 7.83$) y Vue.js ($S = 7.56$),
        siendo seleccionado formalmente por el Comité de Arquitectura
        Tecnológica. Las fortalezas de Angular en seguridad nativa,
        mantenibilidad, cumplimiento regulatorio y --de manera
        determinante-- en la calidad y gobernanza del código asistido por
        IA, resultaron decisivas.

  \item El sistema de indicadores técnicos y regulatorios desarrollado
        permite establecer una línea base mensurable para el seguimiento
        post-implementación de la decisión tecnológica, con mayor
        sensibilidad a las particularidades del contexto bancario que las
        métricas genéricas de proyectos tecnológicos.

  \item La consolidación de herramientas de Inteligencia Artificial
        generativa aplicadas al desarrollo de software no es un factor
        marginal, sino la variable que justifica el propio alcance de esta
        investigación. El análisis de sensibilidad comparado (Tabla
        \ref{tab:4.2}, escenarios S0 vs. S3) mostró que incorporar los
        criterios de calidad y gobernanza de IA \emph{amplía} la ventaja de
        Angular sobre React en 0.10 puntos (de $\Delta S = 0.91$ sin dichos
        criterios a $\Delta S = 1.01$ con ellos), confirmando la Hipótesis
        Específica 6: la disponibilidad de asistentes de código basados en
        IA reduce la relevancia relativa de los criterios organizacionales
        que antes favorecían a frameworks con mayor ecosistema y talento
        disponible (como React), reforzando en cambio la preferencia por
        frameworks tipados y opinionados como Angular, cuya arquitectura
        produce código IA-asistido más consistente, auditable y
        conforme con las exigencias de trazabilidad regulatoria del sector
        bancario. En consecuencia, la emergencia de la IA generativa no
        elimina la necesidad de un proceso de selección tecnológica
        riguroso: reconfigura sus criterios decisivos hacia lo técnico y lo
        regulatorio, exactamente el alcance que da nombre a esta tesis.
\end{enumerate}

\section*{Líneas de investigación futura}

La presente investigación abre al menos tres líneas de trabajo que pueden
enriquecer y extender la metodología propuesta:

\begin{itemize}
  \item \textbf{Actualización de pesos ante la adopción masiva de IA
        generativa.} Aplicar una nueva ronda Delphi con el mismo panel de
        expertos en el horizonte 2027--2028 para medir cómo la
        generalización de asistentes de código (Copilot, Cursor y
        equivalentes) modifica la valoración de los criterios de calidad y
        gobernanza de IA, y determinar si el ranking Angular $>$ React $>$
        Vue.js se mantiene bajo esa nueva distribución de pesos.

  \item \textbf{Extensión a arquitecturas de micro-frontends.} La
        metodología fue diseñada para la selección de un único framework
        base. Su adaptación a escenarios de micro-frontends --donde
        distintos equipos pueden emplear frameworks diferentes dentro de
        una misma plataforma-- requiere incorporar criterios técnicos
        adicionales de interoperabilidad que no fueron contemplados en el
        presente estudio.

  \item \textbf{Validación en instituciones bancarias de mayor escala.} El
        caso de estudio corresponde a una institución mediana con 20
        profesionales de desarrollo. Aplicar la metodología en un banco de
        primer piso con equipos de desarrollo superiores a 200 personas y
        mayor madurez en la adopción de IA generativa permitiría verificar
        si los pesos AHP y el ranking de frameworks se mantienen estables.
\end{itemize}
